\cotefile{156-9}
\scanpage{42}
\mspage{38}
% archive sheet number: 41

y jouant le rôle des $\mathrm{Inf}$ quelconques, contenant $\varnothing$ (plus petit élément) et $\unsure{\mathcal A|A}$ (plus grand élément). On y a une opération $\unsure{\mathrm{cosupp}^{\circ}}$ ou $C$, qui est une anti-involution. Sa donnée \deleted{implique} (p.ex) implique l'existence des $\mathrm{Sup}$ quelconques dans $\Sigma_{\unsure{\mu|M|U}}$ (anti $\unsure{\mathrm{Inf}|\mathrm{Sup}}$) et $C$ échange $\mathrm{Inf}$ et $\mathrm{Sup}$, plus petit et plus grand élément. La relation $\mathrel{|\circ|}$ \deleted{\unsure{se donne par}} \inserted{s'étend} aux éléments de $\Sigma_\mu$, par
\[
S\mathrel{|\circ|}S',\quad \overset{\inserted{\unsure{\mathrm{Not}|\mathrm{def}}}}{\Longleftrightarrow}\quad \forall x\in S,x'\in S',\ \text{on a }x\mathrel{|\circ|}x'.
\]
On a
\[
S\mathrel{|\circ|}S'\Longleftrightarrow S\subset CS'\Longleftrightarrow S'\subset CS.
\]
On a
\[
S\mathrel{|\circ|}S'\Longrightarrow S\cap S'=\varnothing
\]
mais l'inverse n'est pas vrai en général.

% margin note written sideways
% slanted
\margin{Structure d'un \uline{\unsure{quasi-radicié|quasi-radicoïde}}\par
Spatios :\par
a) relation d'ordre\par
b) application $C$\par
\unsure{Alors que} : $1^{\circ})$ $\exists$ plus petit et plus grd élément\par
$2^{\circ})$ $C$ est une anti-involution\par
$3^{\circ})$ \deleted{$\unsure{\forall x\in\Sigma}$} $x\wedge Cx=\varnothing_\Sigma$\par
$(\Longleftrightarrow x\vee Cx=1_\Sigma)$ \unsure{Description|Dans l'option}\par
$\longrightarrow\ \unsure{\mathrm{Figel}\,\mathrm{ord}|\mathrm{Figbord}}(\Sigma)$.}

La donnée d'un magasin \unsure{abstrait|étalon} $\mathcal A$, soit $\unsure{\mathscr F|\mathcal F|\mathbb F},$ (comme ``figures''), \deleted{permet} plus la donnée de $\trianglelefteq$, permet de définir dans $\Sigma_{\mathcal A}$ une partie, formée des \uline{supports constructibles}
